\chapter{Results}
\section{Identifying Emission lines in Hydrogen}
\section{Identifying Absorbtion lines in the Solar Spectra}
\section{Determining the temperature of the Sun}
To determine the temperature of the sun we find the wavelength with maximum intensity, 
in each of the 3 measurements we took.
We assume a constant $\pm 2 \text{ nm}$ uncertainty on the wavelength measurement,
which then becomes an uncertainty in the temperature measurement of $\sigma_T = |\frac{T_\text{max}}{\lambda_\text{max}}\sigma_\lambda|$, when we use Wien's displacement law.
The final estimate for the temperature is then given by a weighted mean of our 3 measurements, and becomes 
\[T = 5564.1 \text{ K} \pm 12.6 \text{ K}\]
\section{Precision of our Spectrometer}